\documentclass[12pt]{article}

\usepackage{amsmath, amssymb, amsthm}
\usepackage{geometry}
\usepackage{hyperref}
\usepackage{booktabs}
\usepackage{array}
\usepackage{xcolor}
\usepackage{natbib}

\geometry{margin=1.1in}

\newtheorem{theorem}{Theorem}
\newtheorem{proposition}{Proposition}
\newtheorem{corollary}{Corollary}
\newtheorem{lemma}{Lemma}
\newtheorem{remark}{Remark}
\newtheorem{definition}{Definition}

\title{Spectral Distortion in Debiased Whittle Estimation\\
Under Structured Satellite Sampling and Missing Data:\\
A Unified Spectral Operator Framework}

\author{Joonwon Lee}
\date{April 2026 (working draft)}

\begin{document}

\maketitle
\tableofcontents
\bigskip

%─────────────────────────────────────────────────────────────────────────────
\section{Introduction}
%─────────────────────────────────────────────────────────────────────────────

The Debiased Whittle (DW) estimator \citep{Sykulski2019} is a widely used
frequency-domain method for estimating the parameters of a stationary
Gaussian process (GP).  Its appeal lies in computational scalability via
the Fast Fourier Transform.  Standard implementations, however, assume
that observations lie on a complete regular grid.

In satellite datasets such as GEMS TCO, this assumption fails in two
distinct, physically motivated ways:
\begin{enumerate}
  \item \textbf{Structured missing data}: GEMS retrievals are missing not
        at random but in an \emph{anisotropic, latitude-structured} pattern
        driven by orbital geometry (East--West orbital stripes, viewing
        angle constraints) and systematic cloud contamination.
        The missingness fraction varies strongly with latitude.
  \item \textbf{Irregular sampling / resampling}: Satellite footprints do
        not fall on the analysis grid.
        In the simulation setup used to validate our framework,
        a high-resolution field is generated and then \emph{mapped to
        actual observation locations} by matching each grid cell to its
        nearest true retrieval location.
        This resampling introduces a spatially structured filtering
        effect---not a simple random jitter.
\end{enumerate}
A common fix is \emph{zero-imputation}: missing or off-grid cells are
filled with zero (after centering) before computing the periodogram.
This paper shows that zero-imputation induces a systematic spectral
distortion that biases DW estimates of the variance $\sigma^2$ downward
and inflates the nugget $\tau^2$, while leaving range parameters
approximately unbiased.

\paragraph{Central theoretical contribution.}
The key insight is that zero-imputation is equivalent to multiplying
the field by a masking process, which induces a \emph{spectral convolution}.
We develop this convolution framework in three layers of increasing
realism for the GEMS setting:
\begin{enumerate}
  \item \textbf{Random i.i.d.\ masking} (Theorem~\ref{thm:general_mask}):
        the idealised tractable case; produces a flat white-noise floor.
  \item \textbf{Deterministic structured masking}
        (Theorem~\ref{thm:deterministic_mask}):
        the latitude-dependent orbital pattern produces an
        \emph{anisotropic}, frequency-dependent distortion.
  \item \textbf{GEMS observation operator}
        (Theorem~\ref{thm:gems_operator}):
        combines the deterministic mask with the resampling kernel
        (nearest-neighbour or cell-averaging) into a single spectral
        operator.
\end{enumerate}
The i.i.d.\ Bernoulli case of prior work is recovered as a special
corollary of Theorem~\ref{thm:general_mask}.

\paragraph{Empirical motivation.}
Using GEMS TCO July data from 2022--2025 (106 analysis days):
\begin{itemize}
  \item The DW GIM standard error for $\sigma^2$ is $0.3960$,
        almost exactly \emph{twice} the Vecchia Matérn SE of $0.1992$.
  \item The DW GIM SE for $\tau^2$ ($0.0404$) is \emph{lower} than
        Vecchia ($0.0620$): the opposite direction from $\sigma^2$.
  \item Score/gradient test with 300 replicates: DW gradient
        $\nabla\ell(\theta_{\mathrm{true}})$ is not negligible
        (mean $L^\infty$ ratio $= 0.73$), while Vecchia Matérn passes
        (ratio $= 0.19$).
\end{itemize}
These findings are precisely consistent with the unified framework.

\paragraph{Outline.}
Section~\ref{sec:background} reviews DW estimation.
Section~\ref{sec:masking_random} develops spectral distortion theory
for random masking.
Section~\ref{sec:masking_det} treats deterministic structured masking,
the latitude-separable GEMS model, and the resampling operator.
Section~\ref{sec:unified} states the unified GEMS operator theorem.
Sections~\ref{sec:pseudo}--\ref{sec:score} derive pseudo-true parameters,
estimation uncertainty, and gradient misspecification.
Sections~\ref{sec:empirical}--\ref{sec:bias} present empirical results
and a bias-correction procedure.


%─────────────────────────────────────────────────────────────────────────────
\section{Background: Debiased Whittle Estimation}
\label{sec:background}
%─────────────────────────────────────────────────────────────────────────────

\subsection{Setup}

Let $\{Y(\mathbf{s})\}$ be a zero-mean stationary GP on a $d$-dimensional
domain with spectral density $f(\boldsymbol{\omega};\theta)$,
$\theta\in\Theta$.  Given $N$ observations on a regular grid with spacing
$\delta$, the discrete Fourier transform (DFT) at frequency
$\boldsymbol{\omega}_j$ is
\[
    d_j \;=\; \frac{1}{\sqrt{N}}\sum_{k=1}^N Y(s_k)\,
                e^{-i\langle\boldsymbol{\omega}_j,\,s_k\rangle},
\]
with periodogram $I_j = |d_j|^2$.
Under Gaussianity, the DFT coefficients are approximately i.i.d.\
$\mathcal{CN}(0, f_j)$ for large $N$, giving the Whittle log-likelihood
\[
    \ell_W(\theta) \;=\; -\sum_j \left[
        \frac{I_j}{f_j(\theta)} + \log f_j(\theta)
    \right].
\]
The debiased variant \citep{Sykulski2019} replaces $f_j(\theta)$ with
$\tilde{f}_j(\theta)$, the expected periodogram of the discretised and
tapered process, correcting for tapering and aliasing.

\subsection{Matérn--nugget spectral model}

We work with the Matérn spectral density with nugget:
\[
  f(\boldsymbol{\omega};\theta)
  \;=\; \sigma^2 f_0(\boldsymbol{\omega};\rho)
        + \frac{\tau^2}{(2\pi)^d},
\]
where $f_0(\cdot;\rho)$ is the unit-variance Matérn density with range
$\rho$, $\sigma^2$ is the partial sill, and $\tau^2$ is the nugget.

\subsection{Zero-imputation}

In standard GEMS TCO implementations (\texttt{generate\_Jvector\_tapered}),
grid cells without valid retrievals are set to zero (after centering by
the observed-cell mean), and the DFT is computed on this zero-imputed field.
We refer to this as \emph{zero-imputation} throughout.


%─────────────────────────────────────────────────────────────────────────────
\section{Spectral Distortion from Random Masking}
\label{sec:masking_random}
%─────────────────────────────────────────────────────────────────────────────

\subsection{Random masking model}
\label{sec:random_mask_model}

Let $Y = \{Y_k\}_{k\in\mathbb{Z}^d}$ be a zero-mean, second-order
stationary field with spectral density $f_Y \in L^1(\mathbb{T}^d)$.
Let $M = \{M_k\}$ be a second-order stationary masking process with mean
$p = \mathbb{E}[M_k]$ and covariance $C_M(h) = \mathrm{Cov}(M_0, M_h)$,
with spectral density $f_M$ of the centered mask $\tilde{M}_k = M_k - p$.
Assume $M \perp Y$.

\begin{theorem}[Spectral distortion under random masking]
\label{thm:general_mask}
Under the above model, $X_k = M_k Y_k$ is second-order stationary with
\[
    C_X(h) \;=\; \bigl[p^2 + C_M(h)\bigr]\,C_Y(h),
\]
and spectral density
\[
\boxed{
    f_X(\boldsymbol{\omega})
    \;=\; p^2\, f_Y(\boldsymbol{\omega})
    \;+\; \frac{1}{(2\pi)^d}
          \bigl(f_M \ast f_Y\bigr)(\boldsymbol{\omega}).
}
\]
\end{theorem}

\begin{proof}
By independence of $M$ and $Y$:
$C_X(h) = \mathbb{E}[M_0 M_h]\,\mathbb{E}[Y_0 Y_h]
= (p^2 + C_M(h))\,C_Y(h)$.
The spectral density follows from the Fourier transform and the
convolution theorem $\mathcal{F}[g_1 \cdot g_2]
= (2\pi)^{-d}(\mathcal{F}[g_1] \ast \mathcal{F}[g_2])$. \qed
\end{proof}

\begin{corollary}[i.i.d.\ Bernoulli masking]
\label{cor:iid_mask}
If $M_k \sim \mathrm{i.i.d.\ Bernoulli}(p)$, then $f_M \equiv p(1-p)/(2\pi)^d$
and
\[
    f_X(\boldsymbol{\omega})
    \;=\; p^2\, f_Y(\boldsymbol{\omega})
    \;+\; \frac{p(1-p)\,\sigma^2_Y}{(2\pi)^d}.
\]
The distortion decomposes into an attenuated signal ($p^2 f_Y$) plus
a flat white-noise floor.
\end{corollary}

\begin{corollary}[Spatially correlated masking induces coloured distortion]
\label{cor:corr_mask}
If $C_M(h) \neq 0$ for $h \neq 0$, the masking-induced distortion is
frequency-dependent (coloured):
\[
    f_X(\boldsymbol{\omega})
    \;=\; p^2 f_Y(\boldsymbol{\omega})
    \;+\; \frac{1}{(2\pi)^d}(f_M \ast f_Y)(\boldsymbol{\omega}).
\]
The noise floor peaks at frequencies where $f_M$ has excess power.
Under spatially clustered cloud cover (low-frequency $f_M$), the
distortion concentrates at low spatial frequencies, contaminating
the signal frequencies most relevant for estimating $\rho$.
\end{corollary}

\begin{remark}[Tapering interaction]
\label{rem:taper}
In practice, the DFT input is $X'_k = W_k M_k Y_k$ with deterministic
taper $W_k$.  The DW debiasing step corrects for the fixed taper through
$\tilde{f}_j(\theta)$, but cannot correct for the random masking $M_k$.
Below we set $W_k \equiv 1$ to isolate masking effects; the taper
extension follows by replacing $f_Y$ with the tapered spectrum in all
theorems.
\end{remark}


%─────────────────────────────────────────────────────────────────────────────
\section{Spectral Distortion from Deterministic Structured Sampling}
\label{sec:masking_det}
%─────────────────────────────────────────────────────────────────────────────

\subsection{Motivation: GEMS TCO sampling geometry}

The random masking model of Section~\ref{sec:masking_random} captures the
stochastic component of day-to-day cloud variability, but it does not
account for the \emph{deterministic} sampling structure imposed by orbital
geometry.
In GEMS TCO data:
\begin{itemize}
  \item Coverage fraction varies systematically with latitude due to
        orbital scan angles and repeat period.
  \item On any given day, the pattern of valid retrievals is a
        near-deterministic stripe pattern overlaid by cloud variability.
  \item In our simulation study, a high-resolution field is generated
        and then mapped to the actual retrieval locations by
        nearest-neighbour assignment---a deterministic resampling
        operator, not a random jitter.
\end{itemize}
This motivates a deterministic masking theorem distinct from
Theorem~\ref{thm:general_mask}.

\subsection{Deterministic structured masking}

Let $m = \{m_k\}_{k\in\mathbb{Z}^d}$ be a \emph{fixed} (deterministic)
binary mask, with $m_k \in \{0, 1\}$.
Define $X_k = m_k Y_k$.
Since $m$ is deterministic, $X$ is generally nonstationary; we
characterise the expected periodogram rather than the power spectrum.

\begin{theorem}[Spectral distortion under deterministic masking]
\label{thm:deterministic_mask}
Let $Y$ be a zero-mean stationary field with spectral density $f_Y$.
Let $m = \{m_k\}$ be a deterministic mask with discrete Fourier
transform $\hat{m}(\boldsymbol{\omega})
= N^{-1/2}\sum_k m_k e^{-i\langle\boldsymbol{\omega},s_k\rangle}$.
Then the expected periodogram of $X_k = m_k Y_k$ satisfies
\[
\boxed{
    \mathbb{E}[I_X(\boldsymbol{\omega})]
    \;=\; \frac{1}{(2\pi)^d}
          \int_{\mathbb{T}^d}
          f_Y(\boldsymbol{\lambda})\,
          \bigl|\hat{m}(\boldsymbol{\omega} - \boldsymbol{\lambda})\bigr|^2
          \,d\boldsymbol{\lambda}.
}
\]
\end{theorem}

\begin{proof}
$\mathbb{E}[I_X(\omega)] = \mathbb{E}[|d_X(\omega)|^2]
= N^{-1}\sum_{k,\ell} m_k m_\ell C_Y(s_k-s_\ell) e^{-i\omega(s_k-s_\ell)}$.
By the Wiener relation, as $N\to\infty$ this converges to
$\frac{1}{(2\pi)^d}\int f_Y(\lambda)|\hat{m}(\omega-\lambda)|^2 d\lambda$.
\qed
\end{proof}

\begin{remark}[Comparison with random masking]
The random masking theorem (Theorem~\ref{thm:general_mask}) involves the
spectral density $f_M$ of the \emph{centered} mask process (an
expectation over randomness), while the deterministic masking theorem
(Theorem~\ref{thm:deterministic_mask}) involves $|\hat{m}|^2$ (the
squared magnitude of the mask's Fourier transform).
Both are spectral convolutions, but with structurally different kernels.
For a random i.i.d.\ Bernoulli mask,
$\mathbb{E}[|\hat{M}(\omega)|^2] \propto f_M(\omega) + p^2 \delta_0(\omega)$,
so the two theorems are consistent in expectation.
\end{remark}

\subsection{Latitude-structured (separable) mask}
\label{sec:lat_mask}

\begin{definition}[Latitude-separable mask]
\label{def:lat_mask}
The GEMS TCO mask is modelled as approximately separable:
\[
    m(s) \;=\; m_{\mathrm{lat}}(s_{\mathrm{lat}})\cdot
                m_{\mathrm{lon}}(s_{\mathrm{lon}}),
\]
where $m_{\mathrm{lat}}$ captures the latitude-dependent orbital coverage
pattern and $m_{\mathrm{lon}}$ captures longitude-dependent cloud effects.
\end{definition}

\begin{corollary}[Anisotropic spectral distortion under separable mask]
\label{cor:anisotropic}
Under the latitude-separable mask of Definition~\ref{def:lat_mask},
\[
    |\hat{m}(\boldsymbol{\omega})|^2
    \;=\; |\hat{m}_{\mathrm{lat}}(\omega_{\mathrm{lat}})|^2\cdot
           |\hat{m}_{\mathrm{lon}}(\omega_{\mathrm{lon}})|^2,
\]
and the spectral distortion in Theorem~\ref{thm:deterministic_mask} is
\emph{anisotropic}: it is concentrated at frequencies where
$|\hat{m}_{\mathrm{lat}}|^2$ or $|\hat{m}_{\mathrm{lon}}|^2$ has excess
power.
If $m_{\mathrm{lat}}$ varies slowly with latitude (wide stripes),
$|\hat{m}_{\mathrm{lat}}|^2$ is concentrated at low latitude frequencies,
causing the distortion to affect the low-frequency spectral content most
relevant for estimating the latitude range parameter $\rho_{\mathrm{lat}}$.
\end{corollary}

\begin{remark}[Implication for range bias]
Corollary~\ref{cor:anisotropic} partially qualifies the ``range is
unbiased'' result of Proposition~\ref{prop:pseudo} (Section~\ref{sec:pseudo}).
Under i.i.d.\ Bernoulli masking, the distortion is isotropic (flat floor)
and does not affect spectral shape.
Under the GEMS latitude-structured mask, the distortion concentrates at
certain directional frequencies, potentially biasing the latitude range
$\rho_{\mathrm{lat}}$ more than the longitude range $\rho_{\mathrm{lon}}$.
Whether this anisotropic bias is significant in practice depends on the
severity of the latitude coverage gradient $m_{\mathrm{lat}}$.
\end{remark}

\subsection{Resampling operator: GEMS simulation setup}
\label{sec:resample}

In the GEMS simulation study, a high-resolution field $Y$ on a fine grid
is generated, and then for each coarse grid cell $k$, the assigned value
is taken from the nearest actual satellite retrieval location:
\[
    X(s_k) \;=\; m_k \sum_i w_{k,i} Y(s_i^{\mathrm{fine}}),
\]
where $w_{k,i}$ are the resampling weights (nearest-neighbour: $w_{k,i}
\in \{0,1\}$, with $\sum_i w_{k,i} = 1$ for observed cells).

\begin{remark}[Jitter vs.\ resampling]
This construction is fundamentally different from the jitter model
$\tilde{s}_k = s_k + \varepsilon_k$ of earlier literature
\citep{Masry1978, Beutler1970}, which assumes a continuous perturbation
of the observation location.
The GEMS resampling operator is a \emph{deterministic, discrete
aggregation}: it selects specific values from the high-resolution field
rather than interpolating a continuously perturbed observation.

In the frequency domain, the resampling operator with weights $\{w_{k,i}\}$
acts as a linear filter with transfer function
$H(\boldsymbol{\omega}) = \sum_i w_{k,i} e^{-i\langle\boldsymbol{\omega},
s_i^{\mathrm{fine}} - s_k\rangle}$, so that
$\mathbb{E}[|d_X(\omega)|^2] \approx f_Y(\omega)|H(\omega)|^2$.
For nearest-neighbour resampling with random footprint location within each
cell,
$|H(\boldsymbol{\omega})|^2 \approx 1$
at low frequencies and decays at high frequencies (spatial smoothing).
For cell averaging,
$|H(\boldsymbol{\omega})|^2 = \mathrm{sinc}^2(\delta\omega/2)$
(boxcar filter), providing a sharper low-pass rolloff.
Unlike Gaussian jitter (which gives $e^{-\sigma^2_\varepsilon\omega^2}$),
the boxcar sinc$^2$ rolloff is sharper and does not decay exponentially.
\end{remark}


%─────────────────────────────────────────────────────────────────────────────
\section{Unified GEMS Observation Operator}
\label{sec:unified}
%─────────────────────────────────────────────────────────────────────────────

We now combine the deterministic latitude-structured mask with the
resampling operator into a single theorem that captures the full
observational process for GEMS TCO.

\begin{theorem}[Spectral distortion under structured satellite sampling]
\label{thm:gems_operator}
Let $Y(s)$ be a stationary random field with spectral density $f_Y$.
Suppose observations are generated as
\[
    X(s_k)
    \;=\; m(s_k)\;\sum_i w_{k,i}\, Y(s_i^{\mathrm{fine}}),
\]
where $m(s_k) \in \{0,1\}$ is the deterministic mask and $H(\boldsymbol{\omega})$
is the transfer function of the resampling operator $\{w_{k,i}\}$.
Then the expected periodogram of $X$ satisfies
\[
\boxed{
    \mathbb{E}[I_X(\boldsymbol{\omega})]
    \;=\;
    \frac{1}{(2\pi)^d}
    \int_{\mathbb{T}^d}
    f_Y(\boldsymbol{\lambda})\,
    |H(\boldsymbol{\lambda})|^2\,
    |\hat{m}(\boldsymbol{\omega} - \boldsymbol{\lambda})|^2\,
    d\boldsymbol{\lambda}.
}
\]
\end{theorem}

\begin{proof}
Write $Z(s_k) = \sum_i w_{k,i} Y(s_i^{\mathrm{fine}})$ (resampled field).
The spectral density of $Z$ is $f_Z(\boldsymbol{\omega}) = f_Y(\boldsymbol{\omega})
|H(\boldsymbol{\omega})|^2$.
The observed process $X_k = m_k Z_k$ is deterministically masked.
Applying Theorem~\ref{thm:deterministic_mask} to $Z$ with spectral density
$f_Z$ gives the result.
\qed
\end{proof}

\begin{remark}[Three-factor decomposition]
Theorem~\ref{thm:gems_operator} decomposes the expected observed
periodogram into three factors:
\begin{itemize}
  \item $f_Y(\boldsymbol{\lambda})$:
        the true spectral content at input frequency $\boldsymbol{\lambda}$;
  \item $|H(\boldsymbol{\lambda})|^2$:
        the resampling filter (low-pass rolloff from nearest-neighbour
        or cell-averaging);
  \item $|\hat{m}(\boldsymbol{\omega}-\boldsymbol{\lambda})|^2$:
        the spectral leakage from the deterministic mask, which mixes
        frequency $\boldsymbol{\lambda}$ into observed frequency
        $\boldsymbol{\omega}$.
\end{itemize}
The total distortion is the integral of $f_Y \cdot |H|^2$ convolved with
$|\hat{m}|^2$.
\end{remark}

\paragraph{GEMS approximation.}
For GEMS TCO with typical parameters, Theorem~\ref{thm:gems_operator}
simplifies as follows.
\begin{enumerate}
  \item \textbf{Resampling filter}: nearest-neighbour resampling on a
        coarse grid (half-degree resolution) means that $|H(\omega)|^2$
        is approximately flat at the analysis resolution; we set
        $|H|^2 \approx 1$ as a first-order approximation.
  \item \textbf{Mask structure}: replacing the deterministic mask $m$
        with its expectation $\bar{m}(s) \approx p$ (mean coverage fraction)
        plus a random centered deviation $\tilde{M}_k$,
        and treating $\tilde{M}_k$ as approximately i.i.d.\ Bernoulli
        deviations with variance $p(1-p)$, Theorem~\ref{thm:gems_operator}
        reduces (in expectation over $\tilde{M}$) to
        Corollary~\ref{cor:iid_mask}:
        \[
          \mathbb{E}[I_X(\boldsymbol{\omega})]
          \;\approx\; p^2 f_Y(\boldsymbol{\omega})
          + \frac{p(1-p)\sigma^2_Y}{(2\pi)^d}.
        \]
\end{enumerate}
This approximation is used in Sections~\ref{sec:pseudo}--\ref{sec:score}
to derive closed-form bias expressions, while
Theorem~\ref{thm:gems_operator} provides the more accurate description
that accounts for latitude-structured coverage anisotropy.


%─────────────────────────────────────────────────────────────────────────────
\section{Pseudo-True Parameters of the Debiased Whittle Estimator}
\label{sec:pseudo}
%─────────────────────────────────────────────────────────────────────────────

The DW estimator minimises the population KL criterion
\[
    \mathcal{K}(\theta)
    \;=\; \int_{\mathbb{T}^d} \left[
        \frac{\mathbb{E}[I_X(\boldsymbol{\omega})]}{f(\boldsymbol{\omega};\theta)}
        + \log f(\boldsymbol{\omega};\theta)
    \right] d\boldsymbol{\omega},
\]
and its limit in probability is the pseudo-true parameter $\theta^*$
\citep{White1982}.

\begin{proposition}[Pseudo-true parameters under i.i.d.\ Bernoulli masking]
\label{prop:pseudo}
Under the GEMS approximation $\mathbb{E}[I_X]
\approx p^2 f_Y + p(1-p)\sigma^2_Y/(2\pi)^d$
and the Matérn--nugget model, the pseudo-true parameter
$\theta^* = (\sigma^{2*}, \rho^*, \tau^{2*})$ satisfies
\begin{align}
  \sigma^{2*} &= p^2\,\sigma^2_{\mathrm{true}},
  \label{eq:sigma_bias} \\[4pt]
  \rho^* &= \rho_{\mathrm{true}},
  \label{eq:range_unbias} \\[4pt]
  \tau^{2*} &= \tau^2_{\mathrm{true}}
               + p(1-p)\,\sigma^2_{\mathrm{true}}.
  \label{eq:nugget_bias}
\end{align}
\end{proposition}

\begin{proof}[Sketch]
The target $\mathbb{E}[I_X]$ lies exactly within the Matérn--nugget model
family with parameters $(\sigma^{2*}, \rho_{\mathrm{true}}, \tau^{2*})$
as given above.  Setting $f(\cdot;\theta) = \mathbb{E}[I_X(\cdot)]$
achieves $\mathcal{K}(\theta) = 1 + $ const (minimum KL), confirming $\theta^*$.
\end{proof}

\begin{remark}[Why range is unbiased: low-frequency argument]
\label{rem:range_unbiased}
Under i.i.d.\ Bernoulli masking, the distortion is isotropic:
$p^2$ scaling of $f_Y$ plus a flat additive constant.
Neither operation changes the spectral shape.
More precisely, for Matérn spectral density
$f_0(\boldsymbol{\omega};\rho) \sim C(1 + \rho^2|\omega|^2)^{-\nu-d/2}$,
at low frequencies the nugget floor $\tau^{2*}/(2\pi)^d$ is negligible
relative to $\sigma^{2*} f_0$, so the shape of $f_0$ is
preserved and identifies $\rho$.
\textit{Range is identified from low-frequency content, where the
masking-induced white-noise floor vanishes asymptotically.}

Under the GEMS deterministic mask (Corollary~\ref{cor:anisotropic}),
this argument partially breaks down: the anisotropic distortion
$|\hat{m}(\boldsymbol{\omega})|^2$ is not flat, and concentrates at
certain directional frequencies, potentially biasing $\rho_{\mathrm{lat}}$.
The severity depends on the gradient of $m_{\mathrm{lat}}$.
\end{remark}

\paragraph{Consequences.}
\begin{enumerate}
  \item \textbf{$\sigma^2$ underestimated}: by factor $p^2 \approx 0.49$
        for $p = 0.70$.
  \item \textbf{$\tau^2$ overestimated}: inflated by $p(1-p)\sigma^2_{\mathrm{true}}$.
  \item \textbf{Range approximately unbiased} under i.i.d.\ masking;
        potentially biased under GEMS latitude structure.
\end{enumerate}


%─────────────────────────────────────────────────────────────────────────────
\section{Effect on Estimation Uncertainty}
\label{sec:fisher}
%─────────────────────────────────────────────────────────────────────────────

The Whittle Fisher information for $\sigma^2$ at pseudo-true $\theta^*$ is
\[
    \mathcal{I}^{\mathrm{DW}}_W(\sigma^2)
    \;=\; \frac{1}{2}
    \int
    \left[
      \frac{p^2 f_0(\boldsymbol{\omega};\rho)}
           {p^2\sigma^2 f_0(\boldsymbol{\omega};\rho) + \tau^{2*}/(2\pi)^d}
    \right]^2 d\boldsymbol{\omega}.
\]
The masking-induced floor $\tau^{2*}/(2\pi)^d$ reduces the contribution of
high-frequency periodogram values to $\mathcal{I}^{\mathrm{DW}}_W(\sigma^2)$,
since there the denominator is dominated by the floor rather than the
signal.

\begin{corollary}[GIM SE inflation for $\sigma^2$]
\label{cor:se}
Under Bernoulli masking with fraction $p$,
\[
  \frac{\mathrm{SE}_{\mathrm{DW}}(\sigma^2)}
       {\mathrm{SE}_{\mathrm{Vecchia}}(\sigma^2)}
  \;\approx\; \frac{1}{p^2}.
\]
For $p = 0.70$: ratio $\approx 2.04$.
\end{corollary}

\begin{proof}[Heuristic]
The DW estimator targets $\sigma^{2*} = p^2\sigma^2$.
The SE of $\hat\sigma^{2*}$ (the estimand) is comparable to the Vecchia
SE for $\sigma^2_{\mathrm{true}}$, since both fit the same spectral family
at their respective parameter scales.
Converting the SE of $\hat\sigma^{2*}$ to the implied SE for $\sigma^2$
introduces a factor $1/p^2$.
\end{proof}

\begin{remark}[Paradoxical SE reduction for $\tau^2$]
The masking floor inflates the apparent nugget signal in $\mathbb{E}[I_X]$.
Intuitively: $\sigma^2$ is estimated from spectral energy, which is
diminished by missing data; $\tau^2$ is estimated from the high-frequency
flat floor, to which zero-imputation contributes a consistent additive
term, paradoxically \emph{increasing} information about the flat component
and thus reducing the SE.
\end{remark}


%─────────────────────────────────────────────────────────────────────────────
\section{Score (Gradient) Test and Model Misspecification}
\label{sec:score}
%─────────────────────────────────────────────────────────────────────────────

\subsection{KL gradient at the true parameters}

For the DW criterion, the population gradient at $\theta$ is
\[
  \mathbb{E}\!\left[\frac{\partial \mathcal{K}}{\partial \theta}\right]
  \;=\; \int
        \left[
          1 - \frac{\mathbb{E}[I_X(\boldsymbol{\omega})]}
                   {f(\boldsymbol{\omega};\theta)}
        \right]
        \frac{\partial \log f(\boldsymbol{\omega};\theta)}{\partial \theta}
        \,d\boldsymbol{\omega}.
\]
At $\theta = \theta_{\mathrm{true}}$, we have $f(\cdot;\theta_{\mathrm{true}}) = f_Y(\cdot)$
but $\mathbb{E}[I_X] = p^2 f_Y + c/(2\pi)^d \neq f_Y$ for $p < 1$.
The dominant bias direction is $\sigma^2$:
\begin{equation}
  \mathbb{E}\!\left[\frac{\partial \mathcal{K}}{\partial \sigma^2}\right]
  \bigg|_{\theta_{\mathrm{true}}}
  \;\propto\; \left(\frac{1}{p^2} - 1\right) > 0,
  \label{eq:sigma_gradient}
\end{equation}
indicating deterministic model misspecification bias (not sampling noise).

\subsection{Score test ratio}

The score/gradient test uses
\[
  r \;=\; \frac{\|\nabla_\theta\ell(\theta_{\mathrm{true}})\|_\infty}
               {\|\nabla_\theta\ell(\theta_{\mathrm{perturbed}})\|_\infty}.
\]
A correctly specified model has $r \ll 1$; large $r$ indicates that
the true parameters are not a stationary point of the expected criterion.

\paragraph{Empirical validation.}

\begin{table}[h]
\centering
\caption{Score test results (mean $L^\infty$ ratio, 300 replicates).}
\begin{tabular}{lccc}
\toprule
Method & Mean ratio (300 iter) & Mean ratio (50 iter) & Interpretation \\
\midrule
Vecchia Matérn (VM) & 0.19 & 0.10 & Pass $\checkmark$ \\
Debiased Whittle (DW) & 0.73 & 0.44 & Fail ($\sigma^2$ bias) \\
Vecchia Cauchy (VC) & 1.47 & 0.73 & Fail (model mismatch) \\
\bottomrule
\end{tabular}
\label{tab:score_test}
\end{table}

VM passes (correctly specified).
VC fails due to model mismatch.
DW fails due to the deterministic $\sigma^2$ gradient bias
in \eqref{eq:sigma_gradient}.


%─────────────────────────────────────────────────────────────────────────────
\section{Empirical Results}
\label{sec:empirical}
%─────────────────────────────────────────────────────────────────────────────

\subsection{Real-data GIM (GEMS TCO July 2022--2025, 106 days)}

\begin{table}[h]
\centering
\caption{GIM standard errors averaged over 106 July days (2022--2025).}
\begin{tabular}{lcccc}
\toprule
Parameter & DW SE & Vecchia-Irr SE & Cauchy-Vecchia SE & DW/Vecchia ratio \\
\midrule
$\sigma^2$ & 0.3960 & 0.1992 & 0.1954 & 1.99 \\
$\tau^2$ (nugget) & 0.0404 & 0.0620 & 0.0879 & 0.65 \\
Range (lon) & --- & --- & --- & --- \\
Range (lat) & --- & --- & --- & --- \\
Range (time) & 0.0726 & 0.0836 & --- & 0.87 \\
\bottomrule
\end{tabular}
\label{tab:gim}
\end{table}

Key observations:
\begin{itemize}
  \item $\mathrm{SE}_{\mathrm{DW}}(\sigma^2)/\mathrm{SE}_{\mathrm{VC}}(\sigma^2)
        = 1.99 \approx 1/p^2 = 2.04$ for $p \approx 0.70$
        (Corollary~\ref{cor:se}).
  \item $\mathrm{SE}_{\mathrm{DW}}(\tau^2) = 0.0404 <
        \mathrm{SE}_{\mathrm{VC}}(\tau^2) = 0.0620$: DW has lower nugget
        uncertainty due to the masking-induced noise floor.
  \item Range SE ratio $\approx 0.87$, close to 1, consistent with
        approximate range unbiasedness under near-i.i.d.\ masking
        (Remark~\ref{rem:range_unbiased}).
\end{itemize}

\subsection{Simulation study (1000 replicates)}

The simulation uses the actual GEMS TCO observation geometry:
a Matérn random field is generated on a fine grid and then mapped to
the actual satellite retrieval locations by nearest-neighbour assignment
(Section~\ref{sec:resample}).
This captures the deterministic resampling structure of
Theorem~\ref{thm:gems_operator}.

\begin{itemize}
  \item DW: RMSRE for $\sigma^2 \approx 0.175$, MdARE $\approx 0.37$.
        Bias consistent with $\sigma^{2*} = p^2\sigma^2_{\mathrm{true}}$
        (Proposition~\ref{prop:pseudo}).
  \item Vecchia Irregular: extreme RMSRE ($\sim 10^{35}$ for $\sigma^2$)
        in some replicates due to near-singular conditioning matrices
        when two actual retrieval locations are nearly co-located.
        A small nugget regularisation (positive diagonal) resolves this.
  \item Vecchia Regular: stable, RMSRE $\approx 0.12$ for all parameters.
\end{itemize}


%─────────────────────────────────────────────────────────────────────────────
\section{Bias Correction for DW Under Missing Data}
\label{sec:bias}
%─────────────────────────────────────────────────────────────────────────────

\subsection{Two-step corrected periodogram}

\begin{definition}[Corrected periodogram]
Given $\hat{p}$ (fraction of non-missing cells) and $\hat{\sigma}^2_Y$,
define
\[
  I_{\mathrm{corr}}(\boldsymbol{\omega})
  \;=\; \frac{I_X(\boldsymbol{\omega})
             - \hat{p}(1-\hat{p})\hat{\sigma}^2_Y / (2\pi)^d}
             {\hat{p}^2}.
\]
\end{definition}

\begin{proposition}[Corrected estimator is approximately unbiased]
Under Corollary~\ref{cor:iid_mask} and $|H|^2 \approx 1$:
$\mathbb{E}[I_{\mathrm{corr}}(\boldsymbol{\omega})] \approx f_Y(\boldsymbol{\omega})$.
\end{proposition}

\paragraph{Iterative estimation.}
Initialise $\hat\sigma^{2(0)}$ from uncorrected DW, then iterate
$\hat\sigma^{2(1)}_Y = \hat\sigma^{2(0)}/\hat{p}^2$ until convergence.

\subsection{Limitations}

\begin{enumerate}
  \item \textbf{Independence}: The correction assumes $m_k \perp Y_k$.
        GEMS cloud cover is positively correlated with aerosol load,
        violating this assumption.
  \item \textbf{Latitude structure / coloured noise floor}:
        The correction removes only a flat floor.
        Under the GEMS latitude-structured mask (Corollary~\ref{cor:anisotropic}
        and Theorem~\ref{thm:gems_operator}), the actual distortion is
        frequency-dependent and anisotropic.
        The residual error after flat-floor correction equals
        $(2\pi)^{-d}\int f_Y(\lambda)[|\hat{m}(\omega-\lambda)|^2
        - p(1-p)]\,d\lambda$,
        which is non-negligible when $|\hat{m}|^2$ departs significantly
        from uniformity.
        A full deconvolution correction requires estimating $|\hat{m}|^2$
        from the binary mask, which is straightforward since $m$ is
        observed.
  \item \textbf{Resampling filter}: For $|H|^2 \not\approx 1$
        (aggressive cell averaging), a frequency-domain deconvolution
        of the filter $|H|^2$ is additionally required.
  \item \textbf{Small $p$}: For $p < 0.3$, the $\hat{p}^2$ denominator
        amplifies noise; spatial methods (Vecchia) are preferred.
\end{enumerate}


%─────────────────────────────────────────────────────────────────────────────
\section{Discussion and Conclusions}
\label{sec:discussion}
%─────────────────────────────────────────────────────────────────────────────

We have developed a unified spectral operator framework that characterises
spectral distortion in the Debiased Whittle estimator under the realistic
satellite observational conditions of GEMS TCO.

The three theorems (random masking, deterministic structured masking,
and the GEMS observation operator) provide a hierarchy of increasingly
accurate descriptions.
The i.i.d.\ Bernoulli approximation (Corollary~\ref{cor:iid_mask})
yields closed-form pseudo-true parameters and SE inflation factors that
match empirical results with quantitative precision.
The deterministic latitude-structured model
(Theorem~\ref{thm:deterministic_mask}, Corollary~\ref{cor:anisotropic},
Theorem~\ref{thm:gems_operator}) reveals that the actual distortion is
\emph{anisotropic} and \emph{frequency-dependent}, with implications for
directional range estimation.

\paragraph{Practical recommendations.}
\begin{itemize}
  \item Use \textbf{Vecchia Matérn} as the primary estimator for $\sigma^2$.
        Raw DW underestimates $\sigma^2$ by $\approx 50\%$.
  \item Use \textbf{DW with two-step correction} as a fast secondary
        estimator; account for latitude structure if
        $|\hat{m}_{\mathrm{lat}}|^2$ departs substantially from flat.
  \item Avoid Vecchia Irregular without nugget regularisation
        (near-duplicate retrieval locations cause numerical instability).
  \item The GEMS simulation design (high-res field $\to$ actual retrieval
        locations) correctly captures the deterministic resampling
        operator $H$; use Theorem~\ref{thm:gems_operator} to characterise
        the resulting periodogram distortion.
\end{itemize}

\paragraph{Future work.}
(a) Estimation of $|\hat{m}|^2$ from binary masks and
    deconvolution-based correction for latitude-structured distortion.
(b) Simulation study quantifying range bias under the GEMS latitude mask.
(c) Extension to the full spatio-temporal DW setting with
    latitude-by-time varying coverage.


%─────────────────────────────────────────────────────────────────────────────
\appendix
\section{Proof of Theorem~\ref{thm:general_mask}}
\label{app:proof1}
%─────────────────────────────────────────────────────────────────────────────

By independence of $M$ and $Y$:
$C_X(h) = \mathbb{E}[M_0 M_h]\,C_Y(h) = (p^2 + C_M(h))\,C_Y(h)$.
Taking the Fourier transform and applying the convolution theorem:
$f_X = p^2 f_Y + (2\pi)^{-d}(f_M \ast f_Y)$.  \qed

\section{Proof of Theorem~\ref{thm:deterministic_mask}}
\label{app:proof2}
%─────────────────────────────────────────────────────────────────────────────

$\mathbb{E}[I_X(\omega)]
= N^{-1}\sum_{k,\ell} m_k m_\ell C_Y(s_k-s_\ell) e^{-i\omega(s_k-s_\ell)}$.
Recognising the double sum as a bilinear form in $C_Y$, and writing
$C_Y(h) = \int f_Y(\lambda) e^{i\lambda h} d\lambda$, the sum converges
as $N\to\infty$ to
$N^{-1}\sum_{k,\ell} m_k m_\ell \int f_Y(\lambda) e^{i(\lambda-\omega)(s_k-s_\ell)} d\lambda
= \frac{1}{(2\pi)^d}\int f_Y(\lambda)|\hat{m}(\omega-\lambda)|^2 d\lambda$.  \qed

\section{DFT Expected Periodogram under i.i.d.\ Bernoulli}
\label{app:dft}
%─────────────────────────────────────────────────────────────────────────────

For i.i.d.\ Bernoulli mask:
$\mathbb{E}[|d^X_j|^2]
= N^{-1}\sum_k \mathbb{E}[M_k^2]\sigma^2_Y
+ N^{-1}\sum_{k\neq\ell} p^2 C_Y(s_k-s_\ell) e^{-i\omega_j(s_k-s_\ell)}$.
Since $\mathbb{E}[M_k^2] = p$ and $p - p^2 = p(1-p)$, this equals
$p^2 f_Y(\omega_j) + p(1-p)\sigma^2_Y/(2\pi)$, matching
Corollary~\ref{cor:iid_mask}.  \qed

\bigskip
\begin{thebibliography}{99}

\bibitem{Sykulski2019}
Sykulski, A.M., Olhede, S.C., Guillaumin, A.P., Lilly, J.M., and Early, J.J.\
(2019).
The debiased Whittle likelihood.
\emph{Biometrika}, \textbf{106}(2), 251--266.

\bibitem{Masry1978}
Masry, E.\ (1978).
Alias-free sampling: an alternative conceptualization and its applications.
\emph{IEEE Transactions on Information Theory}, \textbf{24}(3), 317--324.

\bibitem{Beutler1970}
Beutler, F.J.\ (1970).
Alias-free randomly timed sampling of stochastic processes.
\emph{IEEE Transactions on Information Theory}, \textbf{16}(2), 147--152.

\bibitem{ShapiroSilverman1960}
Shapiro, H.S.\ and Silverman, R.A.\ (1960).
Alias-free sampling of random noise.
\emph{Journal of the Society for Industrial and Applied Mathematics},
\textbf{8}(2), 225--248.

\bibitem{White1982}
White, H.\ (1982).
Maximum likelihood estimation of misspecified models.
\emph{Econometrica}, \textbf{50}(1), 1--25.

\bibitem{Guinness2018}
Guinness, J.\ (2018).
Permutation and grouping methods for sharpening Gaussian process
approximations.
\emph{Technometrics}, \textbf{60}(4), 415--429.

\bibitem{KatzfussGuinness2021}
Katzfuss, M.\ and Guinness, J.\ (2021).
A general framework for Vecchia approximations of Gaussian processes.
\emph{Statistical Science}, \textbf{36}(1), 124--141.

\bibitem{Stein1999}
Stein, M.L.\ (1999).
\emph{Interpolation of Spatial Data: Some Theory for Kriging}.
Springer, New York.

\bibitem{LittleRubin2002}
Little, R.J.A.\ and Rubin, D.B.\ (2002).
\emph{Statistical Analysis with Missing Data}, 2nd ed.
Wiley, Hoboken.

\bibitem{Fuentes2007}
Fuentes, M.\ (2007).
Approximate likelihood for large irregularly spaced spatial data.
\emph{Journal of the American Statistical Association},
\textbf{102}(477), 321--331.

\bibitem{Cressie1993}
Cressie, N.\ (1993).
\emph{Statistics for Spatial Data}, revised ed.
Wiley, New York.

\bibitem{Priestley1981}
Priestley, M.B.\ (1981).
\emph{Spectral Analysis and Time Series}.
Academic Press, London.

\end{thebibliography}

\end{document}
